\chapter{Conclusion}
%
\label{chap:Conclusion}
%
This thesis investigated whether large language models can reliably extract
structured insights from industrial service tickets written primarily in
German.  The investigation was carried out on a 250-ticket benchmark
from Amazone~Werke spanning three ticket categories~(GS, BL, LB) and
fifteen model configurations.

%----------------------------------------------------------------------
\section{Summary of Work}
\label{sec:summary}
%
The work proceeded in three stages.  The first stage was data engineering: a
preprocessing pipeline was built to convert heterogeneous ticket
attachments~(MSG, PDF, DOCX, XLSX) into a unified Markdown representation, and a
rule-based pattern-extraction layer was added to capture structured
fields (document categories, article numbers, machine serial numbers, and RMA
codes) without invoking an LLM.  The second stage was benchmark construction:
250~tickets were selected, annotated by a human expert with ground-truth tags
from a curated set of 28~human-defined tags, and split into three category
subsets.  The third
stage was the LLM evaluation itself, in which fifteen model configurations were
assessed using a two-call extraction and tagging protocol; the results were
measured both with token-level F1 against the human annotations and through an
LLM-as-judge rubric capturing tag usefulness and reasoning coherence.

%----------------------------------------------------------------------
\section{Answers to the Research Questions}
\label{sec:rq-answers}
%
\paragraph{RQ\,1 (Main): Can LLMs be used to reliably extract structured insights
from industrial service tickets?}
The experiments confirm that LLMs are a viable tool for this task,
\emph{provided the prompts are carefully designed}.  Models with sufficiently
large context windows~($\geq$~128k tokens) and moderate temperatures can
process full email threads and produce structured JSON outputs that match
human annotations at an F1 level of up to 0.39 for zero-shot tagging.
While fully replacing human annotation remains a challenge, the models
provide significant value in automated triage and insight synthesis.

\paragraph{RQ\,1.1: Which ticket information can be captured reliably with
rule-based extraction, and which parts need LLMs?}
The hybrid pipeline approach is supported by the evidence.  Structured,
pattern-conformant information such as the document category and the article
number is extracted reliably with regular expressions, reaching an F1 around
0.94 to 1.0 on the manually checked Belastungsanzeige sample.  The machine
serial number, by contrast, is frequently absent from the source document and
could not be treated as a reliable rule-based field.  Narrative content such as
the problem, its cause, and the semantic tags does not follow a fixed pattern
and therefore requires LLM reasoning, which is where rule-based methods fall
short.

\paragraph{RQ\,1.2: Which model produces the most useful and coherent insights
on the benchmark tickets?}
The answer depends on whether proprietary models are permitted.  Across the
fifteen configurations, the proprietary \texttt{gpt-5.4-mini} achieved the
highest token-level F1~(0.394) against the human tags.  Among the locally
hosted open-source models, however, \texttt{gpt-oss:20b}~(GPTOSS) was the
strongest overall: it ranked first in the LLM-as-judge evaluation for both
Problem/Solution quality~(6.09/10) and tag usefulness~(65.6\,\%), despite
ranking only fourth on raw F1~(0.356).  GPTOSS outperformed larger models such
as \texttt{qwen3:14b}, which suggests that instruction tuning and architectural
efficiency matter more than raw parameter count in the 14B to 20B range.

%----------------------------------------------------------------------
%----------------------------------------------------------------------
\section{Future Work}
\label{sec:future-work}
%
Several directions could build on the findings of this thesis.  One of the most
promising is fine-tuning: training a 7B to 14B model on the annotated benchmark
could close the gap with GPTOSS at significantly lower inference cost, making a
locally hosted model competitive with the proprietary alternatives.  This is
closely tied to enlarging the benchmark itself.  Extending it beyond the current
250 tickets to a thousand or more, using model-assisted annotation with a human
in the loop, would improve statistical reliability while keeping the manual
effort manageable.  A complementary step would be a multi-annotator study, in
which a subset of the benchmark is re-annotated by a second expert and the
agreement is quantified with Cohen's~$\kappa$, giving a more robust evaluation
baseline than the single-annotator self-agreement used here.

Beyond strengthening the evaluation, the system could be moved closer to
practice.  Deploying the extraction pipeline as a real-time enrichment service
for incoming tickets, with a confidence threshold below which a human reviewer
is notified, would test the approach under production conditions.  On the
modelling side, prompting in German to match the dominant input language may
improve recall for the smaller models that struggle with cross-lingual
instruction following.  Finally, replacing the free-text generation step with
constrained decoding, such as grammar-based sampling, would guarantee
schema-valid JSON output and remove the need for the post-hoc repair parsing
that the current pipeline relies on.

